\color{unimain}{\subsection{Логика отключения (авто)трансформатора}}
\color{black}

\begin{enumerate}

\item
Ввод/вывод функции осуществляется либо посредством программного ключа \textbf{«ЛО Т(АТ)»}, либо с использованием внешнего сигнала \textit{«ЛО: Вывод»}.

\item
Функциональный блок логики отключения предназначен для подачи сигнала на отключение (авто)трансформатора.

\item
Команда на отключение \textit{«ЛО: Отключение»} формируется при срабатывании логики отключения ошиновки. Команда на запрет АПВ Т(АТ) \textit{«ЛО: Запрет АПВ»} формируется при срабатывании логики запрета АПВ ошиновки.

\item
Уставки логики отключения Т(АТ) представлены в таблице \ref{app_set:lot}. Алгоритм работы логики отключения В1 приведена на рисунке \ref{pict_lo:lot}.

\medskip
\begin{figure}[H]
\centering
\includegraphics[width=0.8\textwidth,height=0.8\textheight,keepaspectratio]{\fbpath/05.02 ДЗО/761.9901 Логика отключения Т/img/lot.pdf}
\captionof{figure}{Функциональный блок логики отключения Т(АТ)}\label{pict_lo:lot}
\end{figure}


\end{enumerate}